% Pick--Poisson determinant: Vandermonde / resultant / discriminant identities.

\begin{theorem}[Pick--Poisson 行列式的 Vandermonde--结果式恒等式]\label{thm:xi-pick-poisson-vandermonde-resultant-identity}
取互异点集 \(A=\{a_1,\dots,a_\kappa\}\subset\mathbb D\)，并令 Pick--Poisson Gram 矩阵 \(\mathsf P\) 由 \eqref{eq:xi-pick-poisson-gram} 给出。
则 \(\det\mathsf P>0\)，并满足乘积分解（命题 \ref{prop:xi-pick-poisson-det-factorization} 的记号）
\[
\det\mathsf P
=
\Bigl(\prod_{j=1}^{\kappa} \frac{1-|a_j|^2}{|1+a_j|^2}\Bigr)\cdot
\Bigl(\prod_{1\le j<k\le \kappa}\rho(a_j,a_k)^2\Bigr),
\qquad
\rho(a,b):=\left|\frac{a-b}{1-\overline b a}\right|.
\]
进一步，令
\[
F(z):=\prod_{j=1}^\kappa (z-a_j)\quad(\text{首一多项式}),\qquad
F^\ast(z):=z^\kappa\,\overline{F(1/\overline z)}=\prod_{j=1}^\kappa (1-\overline a_j z).
\]
则成立纯代数不变量恒等式
\[
\det\mathsf P
=
\frac{\displaystyle \prod_{1\le j<k\le\kappa}|a_j-a_k|^2\ \cdot\ \prod_{j=1}^\kappa (1-|a_j|^2)^2}
{\displaystyle |\Res(F,F^\ast)|\ \cdot\ \prod_{j=1}^\kappa |1+a_j|^2},
\]
以及平方化形式
\[
(\det\mathsf P)^2
=
\frac{\displaystyle |\Disc(F)|\ \cdot\ \prod_{j=1}^\kappa (1-|a_j|^2)^4}
{\displaystyle |\Res(F,F^\ast)|^2\ \cdot\ \prod_{j=1}^\kappa |1+a_j|^4}.
\]
\leanverified{paper\_zeta\_vandermonde\_resultant\_seeds}
\leanverified{paper\_zeta\_vandermonde\_resultant\_package}
\leanverified{paper\_xi\_pick\_poisson\_vandermonde\_resultant\_identity}
\end{theorem}
\begin{proof}
由定义
\[
\prod_{1\le j<k\le\kappa}\rho(a_j,a_k)^2
=
\prod_{1\le j<k\le\kappa}\frac{|a_j-a_k|^2}{|1-\overline a_k a_j|^2}.
\]
另一方面，结果式的根值表述给出
\[
\Res(F,F^\ast)=\prod_{j=1}^\kappa F^\ast(a_j)
=\prod_{j=1}^\kappa\ \prod_{k=1}^\kappa (1-\overline a_k a_j).
\]
取模得
\[
|\Res(F,F^\ast)|
=
\Bigl(\prod_{j=1}^\kappa (1-|a_j|^2)\Bigr)\cdot
\Bigl(\prod_{1\le j<k\le\kappa}|1-\overline a_k a_j|^2\Bigr).
\]
将该式与 \(\rho\) 的展开代入 \(\det\mathsf P\) 的乘积分解并整理，即得第一恒等式。
最后，由首一多项式判别式满足
\(
|\Disc(F)|=\prod_{1\le j<k\le\kappa}|a_j-a_k|^2\cdot \prod_{1\le j<k\le\kappa}|a_j-a_k|^2
=\prod_{j<k}|a_j-a_k|^4
\)，
从而平方化得到第二恒等式。证毕。
\end{proof}

\begin{corollary}[行列式下界推出伪双曲分离的充分条件]\label{cor:xi-pick-poisson-determinant-separation-sufficient}
\leanverified{paper\_xi\_pick\_poisson\_determinant\_separation\_sufficient}
在定理 \ref{thm:xi-pick-poisson-vandermonde-resultant-identity} 的设定下，若进一步满足 \(\max_j|a_j|\le r<1\)，则
\[
\min_{1\le j<k\le\kappa}\rho(a_j,a_k)
\ \ge\
(\det\mathsf P)^{1/2}\,(1-r)^{\kappa}.
\]
\end{corollary}
\begin{proof}
由 \(|a_j|\le r\) 得 \(|1+a_j|\ge 1-|a_j|\ge 1-r\)，故
\(
\frac{1-|a_j|^2}{|1+a_j|^2}\le (1-r)^{-2}
\)。
代入 \(\det\mathsf P=\bigl(\prod_j \frac{1-|a_j|^2}{|1+a_j|^2}\bigr)\bigl(\prod_{j<k}\rho_{jk}^2\bigr)\) 得
\[
\prod_{1\le j<k\le\kappa}\rho(a_j,a_k)^2
\ \ge\
\det\mathsf P\cdot (1-r)^{2\kappa}.
\]
由于每个 \(\rho(a_j,a_k)^2\in(0,1]\)，故其乘积不超过任何单个因子，特别地
\[
\prod_{j<k}\rho(a_j,a_k)^2
\ \le\
\min_{j<k}\rho(a_j,a_k)^2.
\]
合并并取平方根即得结论。证毕。
\end{proof}

\endinput
